\section{Completion of the undifferentiated Stieltjes edge}
\label{edge:sec}

The positive-argument-derivative product theorem in the earlier
\emph{Polygamma Tail Identities} report \cite{tail-polygamma} assumes that both ordinary
derivative orders are positive.  Its research question 9 asks for the
missing cases in which one factor is undifferentiated.  We resolve
that edge.  The mechanism remains the classical Mikol\'as product
integral; the additional result is the joint spectral-pole completion,
its all-index finite algorithm, and its exact endpoint law.  No claim
of historical priority for the classical product kernel is made.
The five incoming packages at the pinned October 11 revision were
also checked for duplication of this statement; their spectral
Stieltjes defects concern different products.

Use the convention
\[
 \zeta(1+u,x)=\frac1u+
       \sum_{m\ge0}\frac{(-1)^m}{m!}\gamma_m(x)u^m,
 \qquad \gamma_0(x)=-\psi(x).
\]
Every finite part below means the constant coefficient of the integral
from $\epsilon$ to $1$ in the original coordinate $x$.  Terms
$\epsilon^a\log^j\epsilon$ with $a<0$ and all positive powers of
$\log\epsilon$ are deleted.  A change of cutoff coordinate is not
part of this convention.

\subsection{The completed first-order edge}

For an integer $l\ge2$, put
\begin{align}
 P_l(u,v)&=\frac{\Gamma(l+u+v)}{\Gamma(1+u)\Gamma(l+v)}
                              \zeta(l+u+v),\\
 E(u)&=\frac{\pi u}{\sin\pi u},\\
 H_l(u,v)&=-\frac{P_l(u,v)E(u)-P_l(0,v)}u
       +(-1)^l\frac{P_l(u,v)E(v)-P_l(u,0)}v.
 \label{edge:H}
\end{align}
Each divided difference in \eqref{edge:H} is holomorphic near
$(0,0)$.  In particular $H_l$ is a genuine two-variable Taylor germ,
not an iterated choice of Laurent constants.

\begin{theorem}[Every undifferentiated-factor jet]
\label{edge:theorem}
For all integers $m,b\ge0$ and $l\ge2$,
\begin{equation}
 \left.\partial_u^m\partial_v^b H_l(u,v)\right|_{u=v=0}
 =(-1)^m\operatorname{FP}\int_0^1
    \gamma_m(x)\left.\partial_t^b\zeta(t,x)\right|_{t=l}\,dx.
 \label{edge:alljets}
\end{equation}
Consequently every finite part
\[
 \operatorname{FP}\int_0^1\gamma_m(x)\gamma_n^{(r)}(x)\,dx,
 \qquad m,n\ge0,\quad r\ge1,
\]
has a finite expression in ordinary zeta values and derivatives at
positive integers.  The superscript $(r)$ is an argument derivative.
\end{theorem}

\begin{proof}
Write $s=1+u$ and $t=l+v$.  The classical product integral continues as
\begin{align}
 F(s,t)&=\pi\frac{\Gamma(s+t-1)}{\Gamma(s)\Gamma(t)}
              \zeta(s+t-1)\{\csc(\pi s)+\csc(\pi t)\}.
 \label{edge:F}
\end{align}
Initially it equals $\int_0^1\zeta(s,x)\zeta(t,x)\,dx$ when
$\Re s<1$, $\Re t<1$, and $\Re(s+t)<1$; see
Shpot--Chaudhary--Paris, Remark 1 \cite{ShpotParis}, or prove it by integrating the
absolutely convergent Hurwitz Fourier series first for
$\Re s,\Re t<0$.  In that initial region the scalar identity
$\int_0^1\zeta(t,x)\,dx=0$ implies
\[
 \int_0^1\left(\zeta(s,x)-\frac1{s-1}\right)\zeta(t,x)\,dx
 =F(s,t).
\]
Both sides have the same meromorphic continuation.

Set $R(s,x)=\zeta(s,x)-1/(s-1)$.  For $s$ near $1$, this function
is holomorphic in $s$, and
\[
 R(s,x)=x^{-s}+A(s,x),\quad
 A(s,x)=\zeta(s,1+x)-\frac1{s-1}.
\]
The function $A$ is jointly holomorphic near $s=1$, $x=0$.
For the second factor use
\[
 \zeta(t,x)=x^{-t}+B(t,x),\quad B(t,x)=\zeta(t,1+x).
\]
In the product the two exponents crossing $-1$ at $(s,t)=(1,l)$ are
precisely
\[
 \zeta(t)x^{-s},\qquad
 c_{l-1}(s)x^{l-1-t},\quad
 c_j(s)=\frac{(-1)^j(s)_j}{j!}\zeta(s+j).
\]
The subtraction of $1/(s-1)$ changes only the constant Taylor
coefficient of $A$; since $l-1\ge1$, the second coefficient is still
$c_{l-1}(s)$ and is holomorphic near $s=1$.
After subtracting sufficiently many further Taylor monomials, the
remainder and all prescribed parameter derivatives are dominated by
$Cx^{-\delta}(1+|\log x|^N)$, with $\delta<1$.
Their integrals therefore form a holomorphic germ.

The analytic integrals of the two resonant monomials are
$-\zeta(t)/(s-1)$ and $-c_{l-1}(s)/(t-l)$, whereas all their
parameter derivatives at the resonance have coordinate finite part
zero: they are multiples of $x^{-1}\log^j x$.
Every nonresonant monomial commutes with parameter differentiation.
Thus the required holomorphic completion is
\[
 F(s,t)+\frac{\zeta(t)}{s-1}+
                         \frac{c_{l-1}(s)}{t-l}.
\]
Now $P_l(0,v)=\zeta(l+v)$ and
$c_{l-1}(1+u)=-(-1)^lP_l(u,0)$.
Substitution into \eqref{edge:F} gives exactly \eqref{edge:H}.
Finally $\partial_u^mR(1+u,x)|_{u=0}=(-1)^m\gamma_m(x)$,
proving \eqref{edge:alljets} with the stated finite-part convention.
\end{proof}

\subsection{A terminating coefficient algorithm}

Define
\[
 p_{r,j}=\left.\frac{d^j}{du^j}(1+u)_r\right|_{u=0},
 \qquad p_{r,j}=0\ \text{if }j>r.
\]
For $r\ge1$ the Hurwitz derivative law gives
\begin{equation}
 \gamma_n^{(r)}(x)=(-1)^{n+r}
       \sum_{j=0}^{\min(n,r)}\binom nj p_{r,j}
            \left.\partial_t^{n-j}\zeta(t,x)\right|_{t=r+1}.
 \label{edge:conversion}
\end{equation}
Write $P_l(u,v)=\sum_{a,b\ge0}P^{(l)}_{ab}u^av^b$ and
$E(u)=\sum_{j\ge0}e_{2j}u^{2j}$, where
\[
 e_0=1,\qquad e_{2j}=2(1-2^{1-2j})\zeta(2j)\quad(j\ge1).
\]
Then
\begin{equation}
 [u^av^b]H_l=
 -\sum_{j=0}^{\lfloor(a+1)/2\rfloor}e_{2j}P^{(l)}_{a+1-2j,b}
 +(-1)^l\sum_{j=0}^{\lfloor(b+1)/2\rfloor}
                        e_{2j}P^{(l)}_{a,b+1-2j}.
 \label{edge:finite}
\end{equation}
Combining \eqref{edge:conversion} and \eqref{edge:finite} evaluates
any requested $m,n,r$ with coefficients of total degree at most
$m+n+1$.  There are no spectral limits left in the algorithm.
The coefficient source can itself be written as
\begin{equation}
 P_l(u,v)=\zeta(l+u+v)
       \prod_{j=1}^{l-1}\left(1+\frac{u}{j+v}\right)
       \exp\left\{
       \sum_{k\ge2}\frac{(-1)^k\zeta(k)}k
                    [(u+v)^k-u^k-v^k]\right\}.
 \label{edge:coefficient-source}
\end{equation}
The product is finite, and only finitely many terms of the exponential
can affect a prescribed coefficient.  This exhibits the claimed
ordinary-zeta closure directly.
In conjunction with the inherited positive-positive sector this proves
a useful uniform statement: whenever the total ordinary derivative
order is positive, every quadratic same-point Stieltjes finite part
reduces to a finite expression involving ordinary positive-integer
zeta values and their order derivatives.  No Stieltjes constant
is needed in those final expressions.  The double-undifferentiated
corner is different and is recorded below for comparison.

For completeness the edge and inherited interior can be written in one
finite dictionary. For $r,t\ge0$ with $r+t>0$, put
$k=r+1$, $l=t+1$, $w=r+t+1$, and define
\begin{align}
 P_{kl}(u,v)&=\frac{\Gamma(w+u+v)}{\Gamma(k+u)\Gamma(l+v)}
                         \zeta(w+u+v),\notag\\
 H_{kl}(u,v)&=(-1)^k\frac{P_{kl}(u,v)E(u)-P_{kl}(0,v)}{u}
       +(-1)^l\frac{P_{kl}(u,v)E(v)-P_{kl}(u,0)}{v}.
       \label{edge:uniform-kernel}
\end{align}
Then the uniform coefficient formula is
\begin{equation}\label{edge:uniform-dictionary}
 \boxed{\displaystyle
 \operatorname{FP}\int_0^1\gamma_m^{(r)}(x)\gamma_n^{(t)}(x)\,dx
 =(-1)^{m+n+r+t}m!n![u^mv^n]
             (1+u)_r(1+v)_tH_{kl}(u,v).}
\end{equation}
When $r=0$ this is the new edge theorem; when $t=0$ use symmetry.
When both are positive it is the inherited positive-positive formula
\cite[Theorem \texttt{hprod:stieltjes-products}]{tail-polygamma}.
Multiplying its Hurwitz Taylor germ by the two rising factorials
proves \eqref{edge:uniform-dictionary}. Only order derivatives at the
single argument $w=r+t+1$ occur. Their coefficients involve rational
numbers and ordinary zeta values from the Gamma expansion. Thus
``single argument'' refers to the differentiated zeta factors, not
to every undifferentiated zeta value in the coefficients.

\subsection{Explicit first-index identities}

Let $H=H_{l-1}$, $H_2=H_{l-1}^{(2)}$, and
$\sigma=(-1)^l$.  The first coefficients are
\begin{align}
 H_l(0,0)&=-H\zeta(l)+(\sigma-1)\zeta'(l),\\
 \partial_u H_l(0,0)
 &=\left[(\sigma-1)\zeta(2)+\left(\frac12-\sigma\right)H_2
                          -\frac12H^2\right]\zeta(l)
       +(\sigma-1)H\zeta'(l)
       +\left(\sigma-\frac12\right)\zeta''(l),
       \label{edge:firstu}\\
 \partial_v H_l(0,0)
 &=\{(\sigma-1)\zeta(2)+H_2\}\zeta(l)
       -H\zeta'(l)+\left(\frac\sigma2-1\right)\zeta''(l).
       \label{edge:firstv}
\end{align}
For example the $l=2$ formulas give the pair
\begin{align}
 \boxed{\operatorname{FP}\int_0^1
       \gamma_1(x)\gamma_0'(x)\,dx
          =\frac12\zeta''(2)-\zeta(2),}
       \label{edge:example10}\\
 \boxed{\operatorname{FP}\int_0^1
       \gamma_0(x)\gamma_1'(x)\,dx
          =-\zeta'(2)-\frac12\zeta''(2).}
       \label{edge:example01}
\end{align}
Neither individual identity follows merely by integrating the
derivative of $\gamma_0\gamma_1$.  Their sum is the endpoint
constant $-\zeta(2)-\zeta'(2)$.
At one further argument derivative,
\begin{equation}
 \operatorname{FP}\int_0^1\gamma_1(x)\gamma_0''(x)\,dx
  =\left(4\zeta(2)-\frac32\right)\zeta(3)
          +6\zeta'(3)+3\zeta''(3).
 \label{edge:example12}
\end{equation}

One higher Stieltjes index already produces a new derivative order:
\begin{equation}
 \operatorname{FP}\int_0^1\gamma_2(x)\gamma_0'(x)\,dx
 =2\zeta'(2)-\zeta''(2)-\frac23\zeta'''(2)
                         +2\zeta(2)\zeta(3).
 \label{edge:example20}
\end{equation}
More generally, put $J_{0,1}^{m,n}=\operatorname{FP}\int_0^1
\gamma_m(x)\gamma_n'(x)\,dx$.  In its finite formula the
coefficient of the highest ordinary-zeta
order derivative $\zeta^{(m+n+1)}(2)$ is exactly
\begin{equation}
 \frac{(-1)^{m+n+1}(m-n)}{(m+1)(n+1)}.
 \label{edge:top}
\end{equation}
This is a coefficient statement for the displayed reduction and does
not assert arithmetic independence of its evaluated constants.
Indeed the highest derivative can only come from
$-P^{(2)}_{m+1,n}+P^{(2)}_{m,n+1}$ with all Gamma and
Pochhammer factors replaced by their constant terms.  Multiplication
by $(-1)^{m+n+1}m!n!$ gives \eqref{edge:top}.

For readers preferring an ordinary convergent integral, put
$g(x)=\gamma_1(1+x)$ and $h(x)=\psi'(1+x)$.
Equation \eqref{edge:example10} is equivalent to
\begin{equation}
 \begin{split}
 \int_0^1\left\{
   \frac{g(x)-\gamma_1-(\zeta(2)+\zeta'(2))x}{x^2}
  +\frac{\log x}{x}\{h(x)-\zeta(2)\}+g(x)h(x)
                  \right\}\,dx
 =\zeta(2)-\frac12\zeta''(2)+\frac14+\gamma_1.
 \end{split}
 \label{edge:ordinary}
\end{equation}
The first apparent singularity is removable and the second is
$O(\log x)$.  Indeed
$\gamma_1(x)=\log x/x+g(x)$ and $\psi'(x)=x^{-2}+h(x)$;
the removed singular monomials have finite-part integral
$-1/4-\gamma_1$, since those proportional to $x^{-1}$ or
$x^{-1}\log x$ have finite part zero.  This proves
\eqref{edge:ordinary} without any numerical limiting prescription.

\subsection{All endpoint constants and compatibility with integration by parts}

For $r,t,m,n\ge0$ write
\[
 J_{r,t}^{m,n}=\operatorname{FP}\int_0^1
                         \gamma_m^{(r)}(x)\gamma_n^{(t)}(x)\,dx.
\]
Set $p_{0,0}=1$ and $p_{0,j}=0$ for $j>0$, consistently with
$(1+u)_0=1$.

\begin{theorem}[Exact endpoint law]
\label{edge:endpoint}
For every $r,t,m,n\ge0$,
\begin{equation}
 \boxed{\begin{aligned}
 J_{r+1,t}^{m,n}+J_{r,t+1}^{m,n}
 ={}&-\frac{(-1)^{m+r}p_{r,m}}{(r+1)!}
                       \gamma_n^{(r+t+1)}(1)\\
    &-\frac{(-1)^{n+t}p_{t,n}}{(t+1)!}
                       \gamma_m^{(r+t+1)}(1).
 \end{aligned}}
 \label{edge:ibp}
\end{equation}
All constants on the right reduce to positive-integer ordinary-zeta
jets by \eqref{edge:conversion}.
\end{theorem}

\begin{proof}
The shift identity gives
$\gamma_m(x)=x^{-1}\log^m x+\gamma_m(1+x)$ exactly.
Its $r$th derivative has singular part
\[
 S_{m,r}(x)=\frac{(-1)^r}{x^{r+1}}
      \sum_{j=0}^{\min(m,r)}(-1)^j\binom mjp_{r,j}\log^{m-j}x.
\]
The coefficient of $x^{-r-1}\log^0x$ is
$(-1)^{m+r}p_{r,m}$; it vanishes if $m>r$.
The smooth factor is $\gamma_m^{(r)}(1+x)$.
The constant term of the product at zero has three contributions:
the product of its two smooth constant terms, and the two products
of a singular logarithm-free coefficient with the matching smooth
Taylor coefficient.  The singular--singular product has power
$x^{-r-t-2}$ and contributes no constant.
The first contribution cancels the endpoint value at $x=1$.
The remaining two are exactly the negative of the right side of
\eqref{edge:ibp}.  Apply the fundamental theorem of calculus on
$[\epsilon,1]$ and take constant coefficients.
\end{proof}

This law provides a useful consistency check on the completed kernels.
It also shows precisely why the first-derivative edge contains information
beyond the old positive-positive sector: integration by parts determines
$J_{0,1}^{m,n}+J_{1,0}^{m,n}$, but not its antisymmetric part.  Once
the first edge has been established, higher edges can also be related
to the positive-positive formulas while retaining the explicit endpoint
terms.  In particular no formal ``integral of a derivative is zero''
rule is appropriate here.

\subsection{The double-undifferentiated corner}

The corner $r=t=0$ is already treated in the manuscript's
coincident-endpoint Stieltjes chapter.  For compatibility, a useful
holomorphic presentation is as follows.  Define
\begin{align}
 A(u)&=u\zeta(1+u),\\
 C(u,v)&=\frac{\Gamma(1+u+v)}{\Gamma(1+u)\Gamma(1+v)}
   \frac{\pi uv\{\csc(\pi u)+\csc(\pi v)\}}{u+v}
                                    A(u+v).
\end{align}
All apparent singularities in $C$ are removable near zero.  In
particular $C(u,0)=A(u)$ and $C(0,v)=A(v)$.
The completed corner is
\begin{equation}
 H_{11}(u,v)=\frac{A(u)+A(v)-1-C(u,v)}{uv}.
 \label{edge:corner}
\end{equation}
Its Taylor coefficients satisfy
$J_{0,0}^{m,n}=(-1)^{m+n}m!n![u^mv^n]H_{11}$.
This expression follows from the same local subtraction, now applied
to both factors $\zeta(1+u,x)-1/u$ and
$\zeta(1+v,x)-1/v$: the integrated spectral product is
$F(1+u,1+v)+1/(uv)$, and the local resonance corrections are
$(\zeta(1+v)-1/v)/u$ and $(\zeta(1+u)-1/u)/v$.
The two pole subtractions must be made jointly.
Equation \eqref{edge:corner} records compatibility with the inherited
corner rather than claiming it as another new result.


\subsection{Independent verification}

The supplied script constructs the local Gamma germ with exact
rational arithmetic and formal ordinary-zeta derivatives.  It checks
all 81 endpoint identities with $0\le m,n,r,t\le2$, compares the
printed low-index identities, and exports their formulas.  This is
an exact finite verification of the implementation, not a replacement
for the analytic proof at all indices.
The numerical suite computes the left sides by integrating their
convergent local Laurent--logarithmic expansions up to a positive
splitting point and then using ordinary quadrature to $1$.
Thus it does not evaluate the product kernel to obtain both sides
of a test.  Its precision, truncation order, actual residuals, and
second splitting-point check are recorded in the JSON output.
These arbitrary-precision comparisons are diagnostics rather than
interval certificates.
